% Volume VI: Adversaries, Platforms, and Automated Judgment
% Part XII. Platforms and Automated Publicity
% Chapter 69: Algorithmic Accusation
% Spine: proof-spines/ch069-spine.md

\chapter{Algorithmic Accusation}
\label{ch:ch069}


\textbf{Algorithmic accusation} begins where systems convert irregularity into suspicion at scale. Risk scores, anomaly detectors, and related tools can legitimately flag cases for review, but only if institutions preserve the distinction between flagging and violation:
\[
\operatorname{Flag}(x)
\centernot\Rightarrow
\operatorname{Violation}(x).
\]
\Definition\ The flag/violation gap fixes the difference between triage and adjudicated wrongdoing.
A flag marks something as worthy of attention; it does not itself establish guilt, breach, or sanctionable wrongdoing.

The chapter's preface insists that automated systems become accusatory when they forget this non-implication. If a pipeline treats statistical deviation as though it were already a proved offense, then suspicion has been collapsed into judgment. The required discipline is therefore procedural as much as technical: algorithmic tools may prioritize review, but they must not silently convert scores, anomalies, or risk rankings into verdicts. When institutions preserve that separation only by demanding ever more proof without funding timely review, evidentiary rigor flips into an unmeetable burden that leaves automated suspicion in force without accountable resolution.
